% ============================================================
%  Capitolo 5 – Sfide, Etica e Prospettive Future
%              nella Patologia Assistita dall'IA
%  University Pathology Course Handbook
%  Included file — do NOT add \documentclass or \begin{document}
% ============================================================

\chapter{Sfide, Etica e Prospettive Future nella Patologia Assistita dall'IA}
\label{ch:5}

% ============================================================
\section{Introduzione}
% ============================================================

I capitoli precedenti di questo manuale hanno introdotto le basi tecniche della patologia digitale, l'architettura dei modelli di apprendimento automatico e i flussi di lavoro pratici attraverso cui l'intelligenza artificiale (IA) viene integrata nella pratica istopatologica. Sarebbe allettante, a questo punto, concludere che la traiettoria dell'IA in patologia sia quella di un progresso lineare e ininterrotto — una storia in cui dataset sempre più grandi e reti neurali sempre più profonde convergono verso un futuro di diagnosi automatizzata e priva di errori. Tale conclusione sarebbe al tempo stesso prematura e intellettualmente irresponsabile. La storia della tecnologia medica è ricca di innovazioni che promettevano trasformazioni ma hanno prodotto complessità, e l'IA non fa eccezione. La riflessione critica sulle limitazioni, le dimensioni etiche e le conseguenze sociali della patologia assistita dall'IA non è una preoccupazione marginale; è una competenza fondamentale per qualsiasi professionista che lavorerà a fianco di questi sistemi.

Questo capitolo affronta le sfide ancora irrisolte nel settore, i quadri etici che devono governare l'implementazione degli strumenti di IA e le questioni filosofiche che emergono quando alle macchine viene chiesto di svolgere compiti precedentemente considerati di esclusiva competenza di esperti umani qualificati. Guarda anche al futuro, esaminando come i ruoli dei patologi, dei biologi e dei tecnici di patologia siano destinati a evolversi in un panorama sempre più plasmato dall'automazione e dal supporto decisionale basato sui dati. Il capitolo è strutturato per passare dalle limitazioni tecniche (Sezioni~\ref{sec:ch5_limitations} e~\ref{sec:ch5_bias}), attraverso l'interpretabilità e l'epistemologia (Sezioni~\ref{sec:ch5_xai} e~\ref{sec:ch5_philosophy}), all'etica e alla governance (Sezione~\ref{sec:ch5_ethics}), e infine ai futuri pratici dell'automazione e della pratica professionale (Sezioni~\ref{sec:ch5_robotics} e~\ref{sec:ch5_technician}).

Un tema centrale in tutto questo capitolo è che l'adozione dell'IA in patologia non è semplicemente un evento tecnico ma sociotecnico — ridisegna le relazioni tra clinici e pazienti, tra istituzioni e regolatori, e tra competenza umana e inferenza algoritmica. Comprendere queste dimensioni è essenziale non solo per coloro che svilupperanno strumenti di IA, ma, forse ancora più importante, per coloro che li utilizzeranno, valideranno e ne saranno responsabili nelle impostazioni cliniche quotidiane. Lo studente post-laurea che completa questo capitolo dovrebbe essere in grado di confrontarsi criticamente con le affermazioni relative all'IA nella letteratura, di identificare potenziali fonti di danno nei sistemi implementati e di contribuire in modo significativo alle discussioni istituzionali sull'adozione responsabile delle tecnologie di patologia digitale \cite{hanna2024}.

% ============================================================
\section{Limitazioni degli Attuali Sistemi di IA in Patologia}
\label{sec:ch5_limitations}
% ============================================================

\subsection{Fallimento della Generalizzazione e Spostamento del Dominio}

Una delle limitazioni più costantemente documentate dei modelli di IA in patologia computazionale è la loro tendenza a ottenere buone prestazioni su dati che assomigliano al loro set di addestramento, ma a degradarsi sostanzialmente quando applicati a dati provenienti da fonti diverse. Questo fenomeno, noto come \emph{spostamento del dominio} o \emph{spostamento della distribuzione}, sorge perché i modelli di apprendimento profondo apprendono associazioni statistiche specifiche per le condizioni in cui i dati di addestramento sono stati acquisiti. Un modello addestrato su immagini di vetrino intero (WSI) digitalizzate con uno scanner Leica Aperio presso un'istituzione può codificare caratteristiche cromatiche specifiche dello scanner, artefatti di compressione e proprietà di messa a fuoco come parte della sua rappresentazione appresa. Quando lo stesso modello viene applicato a vetrini scansionati con un Hamamatsu NanoZoomer presso un'istituzione diversa, queste caratteristiche di basso livello differiscono e le prestazioni del modello possono calare drasticamente — a volte a livelli non migliori del caso \cite{bankhead2022}.

Lo spostamento del dominio non è limitato alle differenze tra scanner. Le variazioni nei protocolli di fissazione del tessuto, nello spessore delle sezioni, nelle procedure di colorazione con ematossilina ed eosina (H\&E), nei mezzi di copertura e persino nell'età dei reagenti possono introdurre differenze sistematiche tra i dataset. Le istituzioni che hanno sviluppato i propri protocolli di colorazione nel corso di decenni possono produrre vetrini con palette cromatiche caratteristiche che sono invisibili all'occhio umano come fonti di bias, ma che vengono prontamente sfruttate dalle reti neurali come caratteristiche discriminative. Ciò significa che un modello può sembrare classificare accuratamente i sottotipi tumorali mentre in realtà si basa su idiosincrasie di colorazione piuttosto che su segnali morfologici genuini. La conseguenza pratica è che i modelli devono essere rigorosamente validati su coorti esterne prima di qualsiasi implementazione clinica, e che la validazione su un singolo sito esterno è insufficiente per stabilire la generalizzabilità attraverso la diversità della pratica di laboratorio nel mondo reale \cite{song2023}.

Affrontare lo spostamento del dominio è diventato un'area di ricerca attiva. Gli algoritmi di normalizzazione della colorazione — che tentano di mappare la distribuzione cromatica di un'immagine sorgente su uno standard di riferimento — possono parzialmente mitigare la variazione di scanner e colorazione, ma introducono i propri artefatti e non risolvono tutte le fonti di spostamento del dominio. Le tecniche di adattamento del dominio, in cui i modelli vengono messi a punto su piccole quantità di dati del dominio target, offrono un'altra via, ma richiedono dati etichettati dal sito target, che possono essere scarsi o costosi da ottenere. L'apprendimento federato, in cui i modelli vengono addestrati in modo collaborativo tra più istituzioni senza condividere i dati grezzi, rappresenta un approccio promettente per costruire modelli più generalizzabili, ma introduce le proprie sfide tecniche e di governance. La lezione fondamentale per il professionista è che le prestazioni riportate di un modello su un dataset benchmark rappresentano un limite superiore, non inferiore, per le sue prestazioni in un nuovo ambiente clinico.

\subsection{Bias del Dataset ed Effetti Batch}

Grandi dataset pubblici come The Cancer Genome Atlas (TCGA) sono stati preziosi nel consentire la ricerca in patologia computazionale, fornendo migliaia di vetrini digitalizzati collegati a dati genomici, clinici e di sopravvivenza. Tuttavia, questi dataset portano bias sistematici che sono sempre più ben caratterizzati. Howard e colleghi hanno dimostrato che i vetrini TCGA contengono firme specifiche del sito — pattern statistici che codificano l'istituzione contribuente piuttosto che le proprietà biologiche del tumore — e che i modelli addestrati su TCGA possono inavvertitamente imparare a classificare per sito piuttosto che per malattia \cite{howard2021}. Questo effetto batch è una forma di confondimento: il modello apprende un'associazione spuria tra l'origine istituzionale e l'esito di interesse, producendo metriche di prestazione gonfiate che non riflettono un apprendimento biologico genuino.

Gli effetti batch nei dataset di patologia derivano dalle stesse fonti dello spostamento del dominio — tipo di scanner, protocollo di colorazione, processazione del tessuto — ma sono particolarmente insidiosi nei dataset multi-sito perché sono correlati con le variabili biologiche di interesse. Ad esempio, se un sito contribuente è specializzato in un particolare sottotipo tumorale e utilizza anche un protocollo di colorazione distintivo, un modello può imparare ad associare quel pattern di colorazione con il sottotipo, raggiungendo un'elevata accuratezza sui dati trattenuti dello stesso sito mentre fallisce sui dati di altri siti. Un design sperimentale attento — inclusa la cross-validazione stratificata per sito e la validazione esterna prospettica — è essenziale per rilevare e mitigare questi effetti, ma tali design non sono universalmente adottati nella letteratura pubblicata, portando a una sovrastima della prontezza degli strumenti di IA per l'uso clinico.

Le conseguenze degli effetti batch si estendono oltre le metriche di prestazione gonfiate. Quando un modello addestrato su un dataset distorto viene implementato in un contesto clinico, i suoi errori possono essere sistematici piuttosto che casuali — può fallire costantemente su vetrini di particolari scanner, particolari protocolli di colorazione o particolari popolazioni di pazienti. Gli errori sistematici sono più pericolosi degli errori casuali perché è meno probabile che vengano rilevati dal monitoraggio della qualità di routine, che in genere si concentra sulle prestazioni medie piuttosto che sulla distribuzione degli errori tra i sottogruppi. Rilevare e caratterizzare gli errori sistematici richiede un'analisi deliberata dei sottogruppi, che dovrebbe essere una componente standard di qualsiasi studio di validazione dell'IA \cite{howard2021}.

\subsection{L'Illusione dell'Oggettività}

Un malinteso persistente nelle discussioni sulla patologia assistita dall'IA è che i metodi digitali e computazionali siano intrinsecamente più oggettivi della valutazione umana. Questa assunzione confonde la quantificazione con l'oggettività e l'automazione con la neutralità. In realtà, ogni fase della pipeline dell'IA — dalla selezione dei casi di addestramento alla scelta del protocollo di annotazione, dalla progettazione della funzione di perdita alla soglia utilizzata per binarizzare una previsione continua — comporta decisioni umane che incorporano valori, assunzioni e potenziali bias nel sistema. L'output di un modello di IA non è una visione da nessun luogo; è il prodotto di una serie di scelte fatte da sviluppatori, annotatori e implementatori, ognuno dei quali porta le proprie prospettive e vincoli \cite{hanna2024}.

L'illusione dell'oggettività è particolarmente consequenziale perché può sopprimere il controllo critico. Quando un patologo non è d'accordo con la diagnosi di un collega, il disaccordo è visibile e può essere discusso. Quando un patologo non è d'accordo con l'output di un sistema di IA, può esserci una pressione istituzionale o psicologica a deferire all'algoritmo, in particolare se il sistema è stato commercializzato come più accurato o coerente della valutazione umana. Questo fenomeno — a volte chiamato \emph{bias da automazione} — è stato documentato nell'aviazione, nella radiologia e in altri domini ad alto rischio, e non c'è motivo di supporre che la patologia ne sia immune. Formare i professionisti a mantenere un impegno critico con gli output dell'IA, piuttosto che un'accettazione passiva, è quindi una componente essenziale di un'implementazione responsabile dell'IA.

Inoltre, le metriche utilizzate per valutare le prestazioni dell'IA sono esse stesse scelte cariche di valori. Accuratezza, sensibilità, specificità e area sotto la curva caratteristica operativa del ricevitore (AUC-ROC) catturano ciascuna aspetti diversi delle prestazioni del modello e possono portare a conclusioni diverse sull'idoneità di un modello per l'uso clinico. Un modello ottimizzato per l'accuratezza complessiva può avere prestazioni scadenti su casi rari ma clinicamente importanti; un modello ottimizzato per la sensibilità può generare un tasso inaccettabile di falsi positivi. La scelta della metrica di valutazione dovrebbe essere guidata dal contesto clinico e dalle conseguenze dei diversi tipi di errore, non dal desiderio di massimizzare un singolo numero a fini di pubblicazione \cite{bankhead2022}.

\subsection{Obsolescenza del Modello e Deriva del Concetto}

I modelli di IA vengono addestrati su dati che riflettono lo stato della pratica patologica in un particolare momento nel tempo. Man mano che la pratica clinica si evolve — attraverso cambiamenti nei criteri diagnostici, l'introduzione di nuovi biomarcatori, aggiornamenti agli schemi di classificazione dell'OMS o cambiamenti nella popolazione di pazienti — le relazioni statistiche codificate in un modello possono diventare obsolete. Questo fenomeno, noto come \emph{deriva del concetto}, significa che un modello che era stato validato come accurato al momento dell'implementazione può gradualmente diventare meno affidabile senza alcuna modifica al modello stesso. A differenza di un libro di testo, che è visibilmente datato dall'anno di pubblicazione, un modello di IA implementato può continuare a generare previsioni sicure molto tempo dopo che il contesto clinico è cambiato in modi che invalidano le sue assunzioni.

La gestione della deriva del concetto richiede un monitoraggio continuo delle prestazioni del modello in fase di implementazione — una pratica nota come \emph{sorveglianza post-commercializzazione} nella terminologia normativa. Questo è tecnicamente impegnativo, richiedendo la raccolta di etichette di verità di base per un campione rappresentativo di casi elaborati dal modello, ed è organizzativamente impegnativo, richiedendo linee chiare di responsabilità per il monitoraggio e per l'attivazione del riaddestramento o del ritiro del modello. Poche istituzioni dispongono attualmente di quadri robusti per questo tipo di governance longitudinale del modello, e i requisiti normativi per la sorveglianza post-commercializzazione dei dispositivi medici basati su IA/ML sono ancora in evoluzione. Lo studente post-laurea dovrebbe essere consapevole che l'implementazione di un modello di IA non è un evento una tantum ma l'inizio di una responsabilità di gestione continua.

\subsection{Il Divario tra Prova di Concetto e Implementazione Clinica}

La letteratura di patologia computazionale contiene migliaia di articoli che dimostrano che i modelli di IA possono eseguire compiti specifici — rilevamento di tumori, classificazione, previsione di biomarcatori — con un'accuratezza paragonabile o superiore a quella di patologi esperti su dataset benchmark curati. Eppure il numero di strumenti di IA che sono stati integrati con successo nei flussi di lavoro clinici di routine rimane piccolo rispetto al volume della ricerca pubblicata. Questo divario tra prova di concetto e implementazione clinica riflette una serie di sfide che sono in gran parte invisibili nelle pubblicazioni accademiche: la necessità di approvazione normativa, il requisito di integrazione con i sistemi informativi di laboratorio (LIS) e le piattaforme di patologia digitale, le esigenze della validazione clinica in contesti prospettici, i costi di implementazione e manutenzione e la gestione del cambiamento organizzativo necessaria per incorporare nuovi strumenti nei flussi di lavoro consolidati \cite{song2023}.

La sola approvazione normativa può richiedere anni e richiede prove di utilità clinica — non solo accuratezza tecnica — nella popolazione di uso previsto. L'integrazione con l'infrastruttura IT esistente è frequentemente sottovalutata come sfida; molti ospedali operano piattaforme LIS legacy che non sono state progettate per interfacciarsi con strumenti di IA, e il lavoro ingegneristico necessario per costruire interfacce robuste e sicure può essere sostanziale. La validazione clinica in contesti prospettici richiede la cooperazione di patologi, clinici e pazienti, e deve essere progettata per rilevare non solo le prestazioni medie ma anche le modalità di fallimento — i casi in cui il modello è sicuramente sbagliato. Lo studente post-laurea che aspira a contribuire alla traduzione clinica degli strumenti di IA dovrebbe essere preparato alla realtà che il percorso da un algoritmo promettente a uno strumento clinico implementato è lungo, costoso e incerto.

\subsection{Requisiti Computazionali e Costi Infrastrutturali}

I modelli di IA all'avanguardia per la patologia computazionale — in particolare quelli basati su grandi modelli fondazionali o vision transformer addestrati su milioni di WSI — richiedono risorse computazionali sostanziali sia per l'addestramento che per l'inferenza. Addestrare un modello di grandi dimensioni da zero può richiedere centinaia di ore GPU su cluster di calcolo ad alte prestazioni, ponendo tale lavoro al di là della portata della maggior parte dei laboratori individuali o delle istituzioni più piccole. Anche l'inferenza — applicare un modello addestrato a nuovi vetrini — può essere computazionalmente impegnativa durante l'elaborazione di WSI gigapixel, richiedendo un'infrastruttura GPU locale o servizi di calcolo basati su cloud, entrambi con costi finanziari.

Questi requisiti infrastrutturali hanno implicazioni per l'equità nell'accesso agli strumenti di IA. I centri medici accademici ben dotati di risorse nei paesi ad alto reddito sono nella posizione migliore per sviluppare, validare e implementare sistemi di IA avanzati, mentre gli ospedali più piccoli, i laboratori comunitari e le istituzioni nei paesi a basso e medio reddito possono mancare dell'infrastruttura computazionale, delle competenze tecniche o delle risorse finanziarie per partecipare alla rivoluzione dell'IA in patologia. Se gli strumenti di IA vengono sviluppati principalmente su dati di istituzioni ben dotate di risorse e implementati principalmente in quegli stessi contesti, rischiano di esacerbare le disuguaglianze esistenti nella qualità diagnostica piuttosto che ridurle. Affrontare questa sfida richiede uno sforzo deliberato per includere istituzioni diverse nei consorzi di sviluppo dell'IA, per sviluppare modelli computazionalmente efficienti adatti a contesti con risorse limitate e per esplorare modelli open-source e di infrastruttura condivisa che riducano il costo di accesso \cite{chen2023}.

% ============================================================
\section{Bias nell'IA: Fonti, Tipologie e Conseguenze}
\label{sec:ch5_bias}
% ============================================================

\subsection{Comprendere il Bias Algoritmico}

Il bias algoritmico si riferisce a errori sistematici negli output di un sistema di IA che producono risultati ingiusti o non equi per particolari gruppi di individui. Nel contesto della patologia assistita dall'IA, il bias può manifestarsi come accuratezza diagnostica differenziale tra sottogruppi di pazienti definiti per età, sesso, etnia, status socioeconomico o origine geografica. È importante capire che il bias algoritmico non è principalmente una proprietà dell'algoritmo stesso, ma dell'intero sistema sociotecnico in cui l'algoritmo è incorporato — inclusi i dati utilizzati per addestrarlo, le etichette assegnate a quei dati, i processi di misurazione che hanno generato i dati e il contesto in cui il modello viene implementato \cite{hanna2024}. Identificare e mitigare il bias richiede quindi attenzione a ciascuno di questi componenti, non solo alle proprietà matematiche del modello.

Le conseguenze del bias algoritmico in patologia non sono astratte. Un modello che sistematicamente ottiene prestazioni inferiori per i pazienti di gruppi demografici sottorappresentati produrrà più errori diagnostici per quei pazienti — tumori mancati, classificazione errata, raccomandazioni di trattamento inappropriate — contribuendo così alle disuguaglianze sanitarie già ben documentate nella pratica patologica convenzionale. Poiché i sistemi di IA possono elaborare grandi volumi di casi rapidamente, un modello distorto implementato su larga scala può causare danni su una scala che sarebbe impossibile per un singolo professionista distorto. Questo effetto di amplificazione rende l'identificazione e la mitigazione del bias nei sistemi di IA una questione di urgente preoccupazione etica e di salute pubblica \cite{seyyedkalantari2021}.

\subsection{Bias dei Dati: Sottorappresentazione dei Gruppi Demografici}

Il bias dei dati sorge quando il dataset di addestramento non rappresenta adeguatamente la diversità della popolazione a cui il modello verrà applicato. In patologia computazionale, questo si manifesta più comunemente come sottorappresentazione di certi gruppi demografici nelle coorti di addestramento. Grandi dataset pubblici come TCGA sono fortemente orientati verso pazienti di centri medici accademici nordamericani, che sono in modo sproporzionato di ascendenza europea e hanno accesso a cure sanitarie di alta qualità. La morfologia tumorale, la composizione dello stroma, i pattern di infiltrazione immunitaria e altre caratteristiche istologiche che i modelli di IA imparano a riconoscere possono variare sistematicamente tra popolazioni con diversi background genetici, esposizioni ambientali e profili di comorbilità. Un modello addestrato prevalentemente su vetrini di un gruppo demografico può quindi avere prestazioni meno accurate quando applicato a pazienti di gruppi sottorappresentati \cite{chen2023}.

Seyyed-Kalantari e colleghi hanno fornito una dimostrazione sorprendente di questo fenomeno nel contesto dell'interpretazione delle radiografie del torace, mostrando che i modelli di IA addestrati su grandi dataset pubblici hanno sistematicamente sottodiagnosticato le condizioni nei pazienti di gruppi emarginati — incluse donne, pazienti neri e pazienti con status socioeconomico inferiore — anche quando l'accuratezza complessiva del modello era elevata \cite{seyyedkalantari2021}. Sebbene questo studio si sia concentrato sulla radiologia piuttosto che sulla patologia, i meccanismi sottostanti — dati di addestramento che non riflettono la diversità della popolazione di pazienti e metriche di valutazione che aggregano le prestazioni tra i gruppi senza esaminare le disparità dei sottogruppi — sono direttamente applicabili alla patologia computazionale. La lezione è che un'elevata accuratezza complessiva è un criterio necessario ma non sufficiente per un sistema di IA equo; l'analisi delle prestazioni dei sottogruppi è essenziale.

\subsection{Bias delle Etichette: Variabilità Inter-Osservatore nelle Annotazioni}

I modelli di IA per la patologia vengono tipicamente addestrati su dataset in cui le etichette di verità di base — tumore presente/assente, grado 1/2/3, positivo/negativo per un biomarcatore — sono state assegnate da patologi umani. Queste etichette non sono infallibili. La variabilità inter-osservatore nella diagnosi patologica è ben documentata in un'ampia gamma di condizioni, dalla classificazione del cancro alla prostata (dove il sistema di Gleason ha subito molteplici revisioni per migliorare la riproducibilità) alla classificazione dei linfomi e alla valutazione dei linfociti infiltranti il tumore. Quando le etichette di addestramento vengono assegnate da un singolo annotatore, o da un piccolo gruppo di annotatori di una particolare istituzione o tradizione formativa, il modello impara a replicare lo stile diagnostico di quegli annotatori — inclusi i loro bias sistematici e le loro idiosincrasie \cite{bankhead2022}.

Il bias delle etichette è particolarmente problematico quando gli annotatori le cui etichette vengono utilizzate per l'addestramento non sono rappresentativi della più ampia comunità di patologi che utilizzeranno il modello. Se un modello viene addestrato su etichette di subspecialisti esperti di un centro di riferimento terziario, potrebbe non generalizzarsi bene alla pratica diagnostica dei patologi generali negli ospedali comunitari, che possono applicare criteri diversi o avere soglie diverse per i casi equivoci. Al contrario, se le etichette di addestramento vengono raccolte da un gran numero di annotatori con diversi livelli di competenza, il modello può apprendere una regressione verso la media che non riesce a catturare i giudizi sfumati degli specialisti esperti. Non esiste una soluzione semplice al bias delle etichette, ma le migliori pratiche includono l'uso di più annotatori indipendenti, l'aggiudicazione esplicita dei disaccordi e la segnalazione delle statistiche di accordo inter-annotatore insieme alle metriche di prestazione del modello.

\subsection{Bias di Misurazione: Variazione di Scanner e Colorazione}

Il bias di misurazione sorge quando il processo di generazione dei dati introduce differenze sistematiche correlate con l'esito di interesse. In patologia digitale, le fonti più importanti di bias di misurazione sono il tipo di scanner e la variazione della colorazione. Come discusso nella Sezione~\ref{sec:ch5_limitations}, diversi scanner producono immagini con diverse caratteristiche cromatiche, risoluzione e proprietà di compressione. Se la distribuzione dei tipi di scanner in un dataset di addestramento è correlata con la distribuzione delle etichette diagnostiche — ad esempio, se i vetrini di uno scanner hanno maggiori probabilità di provenire da un particolare sottotipo tumorale — il modello può imparare a utilizzare le caratteristiche dello scanner come proxy per la diagnosi \cite{howard2021}.

La variazione della colorazione introduce problemi analoghi. La colorazione H\&E non è una procedura standardizzata; il colore preciso dei depositi di ematossilina ed eosina dipende dalla concentrazione e dall'età dei reagenti, dalla durata delle fasi di colorazione, dal pH delle soluzioni e dalla competenza del tecnico. La colorazione immunoistochimica è ancora più variabile, con differenze nel clone dell'anticorpo, nel protocollo di recupero dell'antigene, nel sistema di rilevamento e nel cromogeno che contribuiscono tutti alla variazione nell'intensità e nel pattern di colorazione. I modelli di IA addestrati su immagini immunoistochimiche di un laboratorio possono quindi non generalizzarsi alle immagini di un altro laboratorio, anche quando la biologia sottostante è identica. La standardizzazione dei protocolli di colorazione e l'uso di strumenti di calibrazione cromatica digitale possono ridurre ma non eliminare questa fonte di bias.

\subsection{Bias di Implementazione e Conseguenze per l'Equità Sanitaria}

Il bias di implementazione si riferisce alla discrepanza tra la popolazione su cui un modello è stato sviluppato e validato e la popolazione a cui viene applicato nella pratica. Anche un modello privo di bias rispetto ai suoi dati di addestramento può produrre risultati non equi se viene implementato in un contesto in cui la popolazione di pazienti, il flusso di lavoro clinico o le risorse istituzionali differiscono da quelle del contesto di sviluppo. Ad esempio, un modello validato in un centro accademico ad alto volume con rapido turnover dei vetrini e revisione da parte di patologi esperti può avere prestazioni diverse in un ospedale comunitario dove i vetrini vengono elaborati meno frequentemente, il controllo della qualità è meno rigoroso e il patologo che esamina gli output dell'IA ha meno competenza subspecialistica \cite{zhang2023}.

L'effetto cumulativo del bias dei dati, del bias delle etichette, del bias di misurazione e del bias di implementazione è un rischio che i sistemi di IA svantaggino sistematicamente i pazienti già più vulnerabili agli errori diagnostici — quelli di gruppi demografici sottorappresentati, quelli serviti da istituzioni con risorse insufficienti e quelli con condizioni rare o scarsamente rappresentate nei dataset di addestramento. Affrontare questo rischio richiede un impegno per l'equità come criterio di progettazione esplicito nello sviluppo dell'IA, non semplicemente come un ripensamento. Ciò significa reclutare attivamente coorti di addestramento diverse, condurre analisi delle prestazioni dei sottogruppi come componente standard della valutazione del modello, coinvolgere le comunità che probabilmente saranno interessate dall'implementazione dell'IA e stabilire requisiti normativi per la validazione incentrata sull'equità \cite{chen2023}.

% ============================================================
\section{Interpretabilità e IA Spiegabile}
\label{sec:ch5_xai}
% ============================================================

\subsection{Il Problema della Scatola Nera}

I moderni modelli di apprendimento profondo — in particolare le reti neurali convoluzionali (CNN) e i vision transformer — raggiungono le loro impressionanti prestazioni attraverso l'apprendimento di rappresentazioni altamente complesse e non lineari distribuite su milioni o miliardi di parametri. Questa complessità è precisamente ciò che consente loro di catturare sottili pattern morfologici che possono sfuggire agli osservatori umani, ma li rende anche opachi: in generale non è possibile ispezionare il funzionamento interno di un modello addestrato e capire, in termini interpretabili dall'uomo, perché ha prodotto un particolare output per un particolare input. Questa opacità è comunemente indicata come il \emph{problema della scatola nera}, e ha implicazioni significative per l'implementazione clinica dei sistemi di IA in patologia \cite{amann2020}.

Il problema della scatola nera non è semplicemente un inconveniente tecnico. In medicina clinica, la capacità di spiegare una diagnosi — di articolare il ragionamento che ha portato dalle osservazioni alle conclusioni — è fondamentale per la pratica della medicina e per la fiducia che i pazienti e i clinici ripongono nei processi diagnostici. Un patologo che riporta una diagnosi di carcinoma duttale invasivo può indicare caratteristiche morfologiche specifiche — pleomorfismo nucleare, figure mitotiche, architettura ghiandolare — che giustificano la diagnosi e possono essere esaminate da un collega che cerca un secondo parere. Un modello di IA che produce la stessa diagnosi senza alcuna indicazione di quali caratteristiche dell'immagine abbiano guidato la previsione non può essere esaminato allo stesso modo. Ciò limita la capacità dei patologi di rilevare errori, di imparare dagli output del modello e di assumersi la responsabilità delle diagnosi informate dall'assistenza dell'IA.

\subsection{Metodi di Attribuzione Basati sul Gradiente e Mappe di Attenzione}

Il campo dell'\emph{IA spiegabile} (XAI) ha sviluppato una serie di metodi per generare spiegazioni post-hoc degli output dei modelli di apprendimento profondo. Tra i più ampiamente utilizzati in patologia computazionale vi sono i metodi di attribuzione basati sul gradiente, di cui la Mappatura dell'Attivazione di Classe Ponderata dal Gradiente (Grad-CAM) è il più prominente. Grad-CAM calcola il gradiente dell'output del modello rispetto alle attivazioni di uno strato convoluzionale specificato, e utilizza questi gradienti per produrre una mappa di calore che evidenzia le regioni dell'immagine di input che hanno influenzato più fortemente la previsione. Quando applicato a patch di WSI, le mappe di calore Grad-CAM possono indicare quali regioni del tessuto il modello ha considerato quando ha preso una decisione di classificazione, fornendo una forma di spiegazione spaziale che può essere sovrapposta all'immagine originale per l'ispezione da parte di un patologo \cite{amann2020}.

I modelli basati sull'attenzione — incluse le architetture di apprendimento multi-istanza (MIL) con pooling di attenzione e i vision transformer — offrono una forma correlata di interpretabilità attraverso le mappe di attenzione. In questi modelli, i pesi di attenzione assegnati a diverse patch di immagini o token durante il passaggio in avanti possono essere visualizzati per mostrare quali parti del vetrino hanno contribuito maggiormente alla previsione finale. A differenza di Grad-CAM, che è una spiegazione post-hoc applicata a un modello che non è stato progettato con l'interpretabilità in mente, le mappe di attenzione sono una proprietà intrinseca dell'architettura del modello e possono quindi fornire una rappresentazione più fedele del processo decisionale del modello. Tuttavia, la ricerca ha dimostrato che i pesi di attenzione non sempre corrispondono a caratteristiche interpretabili dall'uomo, e che pesi di attenzione elevati possono essere assegnati a regioni diagnosticamente irrilevanti \cite{bankhead2022}.

\subsection{Valori SHAP e Attribuzione delle Caratteristiche}

SHapley Additive exPlanations (SHAP) è un framework game-teorico per attribuire l'output di un modello alle sue caratteristiche di input, basato sul concetto di valori di Shapley dalla teoria dei giochi cooperativi. Nel contesto della patologia computazionale, i valori SHAP possono essere utilizzati per quantificare il contributo delle singole caratteristiche dell'immagine — o, nei modelli che incorporano metadati clinici, delle singole variabili cliniche — alla previsione di un modello per un caso specifico. I valori SHAP hanno la proprietà desiderabile di soddisfare diversi assiomi matematici di equità (efficienza, simmetria, dummy e additività), rendendoli una base principiata per l'attribuzione delle caratteristiche. Possono essere calcolati per qualsiasi modello, indipendentemente dall'architettura, rendendoli applicabili all'intera gamma di strumenti di IA utilizzati in patologia \cite{amann2020}.

Nonostante il loro fascino teorico, i valori SHAP e altri metodi XAI hanno importanti limitazioni che devono essere comprese dai professionisti. I metodi di spiegazione post-hoc generano approssimazioni del comportamento del modello piuttosto che descrizioni esatte dei calcoli interni del modello; possono essere sensibili alla scelta della distribuzione di base o di riferimento, e possono produrre spiegazioni plausibili ma errate. Inoltre, la risoluzione spaziale delle spiegazioni basate sul gradiente e sull'attenzione è limitata dall'architettura del modello, e le spiegazioni generate a livello di patch potrebbero non tradursi direttamente in interpretazioni a livello di vetrino. Il campo dell'XAI in patologia è attivo e in rapida evoluzione, e i professionisti dovrebbero approcciare i metodi di spiegazione attuali come strumenti utili ma imperfetti piuttosto che come resoconti definitivi del ragionamento del modello.

\subsection{Spiegabilità, Fiducia Clinica e Approvazione Normativa}

L'importanza della spiegabilità nella patologia assistita dall'IA si estende oltre il dominio tecnico nei regni della fiducia clinica, della responsabilità professionale e della conformità normativa. I clinici sono più propensi a fidarsi e utilizzare appropriatamente gli strumenti di IA quando possono comprendere la base delle raccomandazioni dello strumento, e meno propensi a impegnarsi in un bias da automazione acritico. I pazienti hanno un legittimo interesse a capire come l'IA ha contribuito alla loro diagnosi, in particolare quando la diagnosi assistita dall'IA viene utilizzata per informare decisioni di trattamento consequenziali. Gli organi normativi in più giurisdizioni hanno iniziato a richiedere che i dispositivi medici basati su IA/ML forniscano una qualche forma di spiegazione per i loro output, riflettendo un crescente consenso che l'opacità è incompatibile con un'IA clinica sicura e responsabile \cite{zhang2023}.

Le linee guida di reporting SPIRIT-AI e CONSORT-AI, sviluppate da Cruz Rivera e colleghi, forniscono quadri standardizzati per riportare il design e i risultati degli studi clinici che coinvolgono interventi di IA \cite{cruzirivera2020}. SPIRIT-AI (Standard Protocol Items: Recommendations for Interventional Trials — Artificial Intelligence) estende le linee guida SPIRIT per i protocolli degli studi clinici per affrontare considerazioni specifiche dell'IA, inclusa la specifica dell'intervento di IA, la gestione degli aggiornamenti del modello durante lo studio e la segnalazione delle prestazioni tra i sottogruppi. CONSORT-AI (Consolidated Standards of Reporting Trials — Artificial Intelligence) estende le linee guida CONSORT per la segnalazione degli studi per richiedere la divulgazione di informazioni specifiche dell'IA, inclusa la versione del modello utilizzato, la disponibilità del modello per la validazione indipendente e la natura di qualsiasi supervisione umana degli output dell'IA. L'aderenza a queste linee guida è sempre più attesa dalle riviste ad alto impatto e dagli organi normativi, e la familiarità con esse è una competenza importante per i professionisti coinvolti nella ricerca o nella valutazione dell'IA.

% ============================================================
\section{Riflessioni Filosofiche ed Epistemologiche}
\label{sec:ch5_philosophy}
% ============================================================

\subsection{Cosa Significa per una Macchina ``Diagnosticare''?}

La parola \emph{diagnosi} deriva dal greco \emph{diagign\={o}skein} — distinguere, discernere — e nel suo uso clinico porta connotazioni che si estendono ben oltre l'identificazione di un pattern. Una diagnosi è un'affermazione sulla natura della condizione di un paziente, fatta da un agente responsabile che può essere ritenuto responsabile dell'affermazione, comunicata al paziente e al suo team clinico, e utilizzata come base per decisioni consequenziali sul trattamento e sulla prognosi. Quando diciamo che un sistema di IA ha ``diagnosticato'' un tumore, stiamo usando la parola in un senso che elimina la maggior parte di queste connotazioni: il sistema ha identificato un pattern in un'immagine che corrisponde alla firma statistica di una particolare classe nei suoi dati di addestramento. Se questo costituisca una diagnosi in qualsiasi senso significativo è una questione che merita un attento esame filosofico \cite{hanna2024}.

La distinzione tra riconoscimento dei pattern e comprensione è centrale per questa domanda. I modelli di apprendimento profondo sono, nella loro essenza, sofisticati sistemi di corrispondenza dei pattern: imparano ad associare le rappresentazioni di input con le etichette di output ottimizzando una funzione obiettivo matematica su un grande dataset. Non possiedono credenze, intenzioni o comprensione in alcun senso filosoficamente robusto. Quando una CNN classifica correttamente un vetrino come mostrante carcinoma invasivo, lo fa perché le proprietà statistiche dell'immagine rientrano nella regione dello spazio delle caratteristiche associata a quell'etichetta nei dati di addestramento — non perché il modello capisca cos'è il cancro, perché insorge, come colpisce il paziente o cosa significa la diagnosi per la vita del paziente. Questa distinzione è importante perché la comprensione, nel senso umano, è ciò che consente a un patologo di gestire casi nuovi, di riconoscere quando un caso esula dalla portata della propria esperienza e di comunicare l'incertezza e la sfumatura di una diagnosi ai colleghi clinici \cite{song2023}.

\subsection{Conoscenza Tacita e i Limiti della Formalizzazione}

Il filosofo Michael Polanyi ha introdotto il concetto di \emph{conoscenza tacita} per descrivere la dimensione dell'expertise umana che non può essere completamente articolata in regole o procedure esplicite — la conoscenza che è incorporata nella pratica qualificata e che viene trasmessa attraverso l'apprendistato, l'osservazione e l'esperienza piuttosto che attraverso l'istruzione formale. La diagnosi patologica è ricca di conoscenza tacita: il patologo esperto sviluppa un senso intuitivo di come appare un vetrino ``normale'', una capacità di rilevare sottili deviazioni dalla normalità che resistono a una descrizione precisa, e una capacità di integrare le osservazioni morfologiche con il contesto clinico, la storia del paziente e l'esperienza precedente in modi che non sono facilmente formalizzabili. Questa dimensione tacita dell'expertise patologica è precisamente ciò che è più difficile da catturare in un dataset di addestramento e più resistente alla formalizzazione algoritmica.

Le implicazioni di questa osservazione per l'IA in patologia sono significative. I modelli di IA possono essere addestrati per replicare i componenti espliciti e articolabili della diagnosi patologica — l'identificazione di caratteristiche morfologiche specifiche, l'applicazione di criteri di classificazione definiti — ma non possono, al momento, replicare la dimensione tacita della pratica esperta. Ciò significa che gli strumenti di IA sono probabilmente più affidabili in compiti ben definiti e standardizzati dove i criteri diagnostici sono espliciti e i dati di addestramento sono rappresentativi, e meno affidabili in casi nuovi, ambigui o contestualmente complessi — precisamente i casi in cui il giudizio umano esperto è più prezioso. Riconoscere questa asimmetria è essenziale per progettare appropriati quadri di collaborazione uomo-IA che sfruttino i punti di forza sia dell'intelligenza umana che di quella artificiale \cite{amann2020}.

\subsection{Potenziamento Versus Sostituzione: Prospettive Filosofiche}

La questione se l'IA sostituirà o potenzierà il patologo è stata ampiamente dibattuta nella letteratura e nella stampa popolare, spesso generando più calore che luce. Da una prospettiva filosofica, la domanda non è semplicemente empirica — non è semplicemente una questione di se l'IA possa raggiungere un'accuratezza sufficiente — ma normativa: riguarda che tipo di pratica diagnostica vogliamo, quali valori desideriamo incorporare nei nostri sistemi sanitari e quale ruolo crediamo che l'expertise umana e la responsabilità umana debbano svolgere nel processo decisionale medico. Queste sono domande a cui non si può rispondere solo con le metriche di prestazione del benchmark \cite{hanna2024}.

Diversi quadri filosofici sono rilevanti per questo dibattito. Il concetto di \emph{cognizione distribuita}, sviluppato da Edwin Hutchins, suggerisce che i compiti cognitivi non vengono eseguiti da menti individuali in isolamento ma sono distribuiti attraverso reti di persone, strumenti e artefatti. Da questo punto di vista, il sistema patologo-IA è un sistema cognitivo in cui le capacità umane e quelle della macchina sono complementari e mutuamente costitutive, piuttosto che competitive. Il concetto di IA \emph{human-in-the-loop} formalizza questa intuizione in termini ingegneristici, specificando che la supervisione e l'intervento umano dovrebbero essere mantenuti nei punti decisionali critici in qualsiasi flusso di lavoro assistito dall'IA. Il concetto di \emph{controllo umano significativo}, sviluppato nel contesto dei sistemi d'arma autonomi ma sempre più applicato all'IA medica, va oltre, richiedendo non solo che un essere umano sia presente nel ciclo ma che l'essere umano abbia una comprensione sufficiente degli output del sistema di IA per esercitare un controllo genuino e informato sulla decisione finale \cite{zhang2023}.

\subsection{L'Epistemologia della Diagnosi Assistita dall'IA}

Da una prospettiva epistemologica, la diagnosi assistita dall'IA solleva domande sulla natura e la giustificazione della conoscenza diagnostica. Quando un patologo accetta la raccomandazione di un sistema di IA, su quale base è giustificata tale accettazione? Se il patologo non può ispezionare il ragionamento del modello — perché il modello è una scatola nera — allora l'accettazione si basa sulla fiducia nel track record del sistema piuttosto che sulla comprensione del suo ragionamento. Questo non è necessariamente problematico; i clinici si affidano abitualmente a test di laboratorio e modalità di imaging i cui meccanismi sottostanti non comprendono pienamente. Ma significa che lo status epistemologico della diagnosi assistita dall'IA è diverso da quello della diagnosi patologica tradizionale, e che le condizioni in cui tale fiducia è giustificata — validazione rigorosa, monitoraggio continuo, segnalazione trasparente delle prestazioni — devono essere chiaramente specificate e applicate \cite{cruzirivera2020}.

Le sfide epistemologiche della diagnosi assistita dall'IA sono aggravate dal fenomeno dell'\emph{ingiustizia epistemica}, descritto dalla filosofa Miranda Fricker, che si verifica quando una persona viene danneggiata nella sua capacità di conoscitore. Nel contesto della patologia assistita dall'IA, l'ingiustizia epistemica potrebbe sorgere se i pazienti di gruppi emarginati ricevono diagnosi meno accurate perché i sistemi di IA addestrati su dati non rappresentativi sono meno affidabili per i loro casi, o se i patologi di istituzioni meno prestigiose vengono sistematicamente esclusi dallo sviluppo di strumenti di IA che plasmeranno la pratica diagnostica a livello globale. Prestare attenzione a queste dimensioni della giustizia epistemica fa parte della più ampia responsabilità etica di coloro che sviluppano e implementano sistemi di IA in sanità.

% ============================================================
\section{Etica nella Diagnosi Assistita dall'IA}
\label{sec:ch5_ethics}
% ============================================================

\subsection{Consenso Informato e Diritti del Paziente}

Il principio del consenso informato — che i pazienti hanno il diritto di comprendere e acconsentire alle procedure utilizzate nella loro cura — si applica alla diagnosi assistita dall'IA come a qualsiasi altro intervento medico. Tuttavia, l'applicazione di questo principio all'IA solleva nuove domande. Quali informazioni devono essere divulgate a un paziente affinché il suo consenso alla diagnosi assistita dall'IA sia genuinamente informato? Come minimo, i pazienti dovrebbero essere informati che gli strumenti di IA vengono utilizzati nell'analisi dei loro campioni, qual è lo scopo di tali strumenti e quali sono le limitazioni e i rischi noti degli strumenti. Dovrebbero anche essere informati del loro diritto di richiedere una revisione solo umana se hanno preoccupazioni riguardo al coinvolgimento dell'IA nella loro cura \cite{zhang2023}.

In pratica, l'implementazione del consenso informato per la diagnosi assistita dall'IA è complicata da diversi fattori. Gli strumenti di IA possono essere utilizzati in più fasi del flusso di lavoro diagnostico — nel controllo della qualità, nello screening primario, nella quantificazione dei biomarcatori — e potrebbe non essere pratico ottenere un consenso specifico per ogni applicazione. Il linguaggio utilizzato per descrivere gli strumenti di IA ai pazienti deve essere accessibile e non tecnico, il che è impegnativo data la complessità dei sistemi sottostanti. Inoltre, il concetto di consenso è complicato dal fatto che i modelli di IA vengono addestrati su dati di pazienti passati che potrebbero non aver acconsentito all'uso dei loro dati per lo sviluppo dell'IA. L'uso retrospettivo dei dati dei pazienti per l'addestramento dell'IA è disciplinato dalla legislazione sulla protezione dei dati e dai quadri etici della ricerca, ma i confini dell'uso consentito non sono sempre chiari e le pratiche variano tra le giurisdizioni \cite{hanna2024}.

\subsection{Privacy dei Dati e Proprietà Intellettuale}

Le immagini di vetrino intero del tessuto del paziente sono dati sanitari personali e sono soggette alle protezioni previste dalla legislazione sulla protezione dei dati, incluso il Regolamento Generale sulla Protezione dei Dati (GDPR) nell'Unione Europea e la legislazione equivalente in altre giurisdizioni. L'uso delle WSI dei pazienti per l'addestramento dell'IA richiede basi giuridiche appropriate — tipicamente il consenso esplicito o una determinazione di interesse legittimo — e deve rispettare i requisiti di minimizzazione dei dati, limitazione delle finalità e sicurezza. La de-identificazione delle WSI è tecnicamente impegnativa perché le immagini possono contenere informazioni — come i pattern di morfologia del tessuto che sono unici per un individuo — che potrebbero in linea di principio essere utilizzate per re-identificare il paziente, anche in assenza di identificatori espliciti \cite{zhang2023}.

Lo status di proprietà intellettuale (PI) dei modelli di IA addestrati sui dati dei pazienti è un'area del diritto complessa e in evoluzione. Quando un ospedale o un istituto di ricerca addestra un modello di IA utilizzando i dati dei pazienti, chi possiede il modello risultante? L'istituzione che ha fornito i dati? Gli sviluppatori che hanno progettato e addestrato il modello? I pazienti i cui dati sono stati utilizzati? Gli attuali quadri di PI non sono stati progettati tenendo conto dell'IA, e c'è una significativa incertezza giuridica sulla proprietà degli output generati dall'IA e dei modelli che li producono. Questa incertezza ha conseguenze pratiche: influisce sulla capacità delle istituzioni di commercializzare strumenti di IA sviluppati utilizzando i dati dei pazienti, sulla capacità dei pazienti di beneficiare del valore commerciale dei loro dati e sulla capacità dei regolatori di richiedere l'accesso ai pesi del modello per la valutazione della sicurezza \cite{hanna2024}.

\subsection{Responsabilità e Quadri Normativi}

Quando un sistema di IA contribuisce a un errore diagnostico — un tumore mancato, un grado errato, un risultato falso positivo che porta a un trattamento non necessario — chi è responsabile? Il patologo che ha esaminato l'output dell'IA e lo ha accettato? L'istituzione che ha implementato il sistema di IA? L'azienda che ha sviluppato e venduto il sistema? L'organo normativo che lo ha approvato? Queste domande di responsabilità non sono meramente teoriche; hanno implicazioni dirette per la responsabilità legale dei professionisti e delle istituzioni, per gli incentivi degli sviluppatori di IA e per la progettazione dei quadri normativi \cite{elendu2023}.

Il panorama normativo per i dispositivi medici basati su IA/ML si sta evolvendo rapidamente. Negli Stati Uniti, la Food and Drug Administration (FDA) ha sviluppato un quadro per la regolamentazione del Software come Dispositivo Medico (SaMD) basato su IA/ML, che distingue tra algoritmi \emph{bloccati} (i cui parametri non cambiano dopo l'implementazione) e algoritmi \emph{adattativi} (che continuano ad apprendere da nuovi dati dopo l'implementazione). Gli algoritmi adattativi presentano particolari sfide normative perché la versione dell'algoritmo che è stata validata può differire dalla versione in uso in un dato momento. La FDA ha proposto un quadro di \emph{piano di controllo delle modifiche predeterminato}, in cui gli sviluppatori specificano in anticipo i tipi di modifiche che possono essere apportate a un modello e le procedure di validazione che verranno applicate prima che tali modifiche vengano implementate \cite{zhang2023}.

Nell'Unione Europea, il Regolamento UE sull'IA — adottato nel 2024 — stabilisce un quadro normativo basato sul rischio per i sistemi di IA, classificando le applicazioni di IA in sanità come \emph{ad alto rischio} e sottoponendole a requisiti di trasparenza, supervisione umana, robustezza e accuratezza. I sistemi di IA utilizzati per decisioni diagnostiche o terapeutiche devono essere registrati in un database pubblico, devono essere sottoposti a valutazione di conformità prima dell'implementazione e devono essere accompagnati da documentazione tecnica che consenta il controllo normativo. Il Regolamento UE sull'IA stabilisce anche requisiti per il monitoraggio post-commercializzazione e per la segnalazione di incidenti gravi che coinvolgono sistemi di IA. La conformità a questi requisiti imporrà obblighi significativi agli sviluppatori di IA e alle istituzioni che li implementano, e richiederà lo sviluppo di nuove strutture di governance e capacità tecniche \cite{zhang2023}.

\subsection{Responsabilità Professionale e Pratica Etica}

Al di là dei requisiti formali della legge e della regolamentazione, l'implementazione dell'IA in patologia solleva questioni di responsabilità professionale e pratica etica. I patologi e i biologi hanno obblighi professionali — sanciti nei codici di condotta emanati da organismi come il Royal College of Pathologists e l'Institute of Biomedical Science — di mantenere la qualità e l'accuratezza del loro lavoro diagnostico, di agire nel migliore interesse dei pazienti e di impegnarsi nello sviluppo professionale continuo. Questi obblighi non diminuiscono quando gli strumenti di IA vengono introdotti nel flusso di lavoro; se mai, si intensificano, perché il professionista deve ora assumersi la responsabilità non solo dei propri giudizi ma anche dell'uso appropriato e della valutazione critica degli output dell'IA \cite{hanna2024}.

La responsabilità professionale nel contesto della patologia assistita dall'IA include l'obbligo di comprendere le capacità e le limitazioni degli strumenti di IA utilizzati, di mantenere le competenze necessarie per rilevare gli errori dell'IA, di segnalare le preoccupazioni sulle prestazioni dell'IA attraverso i canali appropriati e di difendere i pazienti i cui interessi possono essere compromessi da sistemi di IA distorti, inaccurati o implementati in modo inappropriato. Include anche l'obbligo di confrontarsi con le più ampie dimensioni etiche dell'IA in sanità — di partecipare alle discussioni istituzionali e professionali sull'adozione responsabile dell'IA, di contribuire allo sviluppo di linee guida etiche e quadri normativi e di garantire che le voci dei pazienti e delle comunità sottorappresentate siano ascoltate in queste discussioni \cite{elendu2023}.

% ============================================================
\section{Robotica, Automazione e Sistemi Autonomi}
\label{sec:ch5_robotics}
% ============================================================

\subsection{Dissezione Robotica e Processazione Automatizzata del Tessuto}

L'automazione dei flussi di lavoro di patologia si estende oltre l'analisi digitale dei vetrini per comprendere la manipolazione fisica dei campioni di tessuto. I sistemi di dissezione robotica — che utilizzano bracci robotici, visione artificiale e protocolli di taglio guidati dall'IA per eseguire l'esame macroscopico e la sezionatura dei campioni chirurgici — sono in fase di sviluppo attivo e sono stati dimostrati in contesti di ricerca. Questi sistemi mirano a standardizzare il processo di dissezione, ridurre la variabilità introdotta dalle differenze nella tecnica del tecnico e migliorare la riproducibilità del campionamento del tessuto. Hanno anche il potenziale per ridurre le richieste fisiche sui tecnici di patologia, che eseguono compiti di taglio ripetitivi che comportano rischi di lesioni muscoloscheletriche e esposizione accidentale a oggetti taglienti \cite{elendu2023}.

Le piattaforme di colorazione automatizzata sono in uso routinario nei laboratori di patologia da molti anni, eseguendo la colorazione H\&E, l'immunoistochimica e l'ibridazione in situ con un grado di standardizzazione e throughput che sarebbe impossibile raggiungere manualmente. Sviluppi più recenti includono sistemi completamente integrati di processazione e colorazione del tessuto che possono portare una cassetta di tessuto dalla fissazione attraverso l'embedding, la sezionatura, la colorazione e la copertura con un intervento umano minimo. I sistemi robotici di gestione dei vetrini — che trasportano i vetrini tra le stazioni di processazione, caricano i vetrini sugli scanner e gestiscono lo stoccaggio e il recupero dei vetrini — sono sempre più comuni nei laboratori di patologia digitale ad alto volume. Insieme, queste tecnologie stanno creando una visione del laboratorio di patologia come un ambiente altamente automatizzato e gestito roboticamente in cui l'intervento umano è riservato al controllo della qualità, alla gestione delle eccezioni e al processo decisionale diagnostico \cite{song2023}.

\subsection{Diagnosi Autonoma tramite IA: Concetto e Controversia}

Il concetto di \emph{diagnosi autonoma tramite IA} — in cui un sistema di IA produce un referto diagnostico senza alcuna revisione umana — rappresenta il punto finale logico della traiettoria di automazione in patologia. I sostenitori sostengono che la diagnosi autonoma tramite IA potrebbe aumentare drasticamente il throughput diagnostico, ridurre i tempi di risposta ed estendere l'accesso a diagnosi di alta qualità in contesti in cui l'expertise del patologo è scarsa. Nei paesi a basso e medio reddito, dove il rapporto tra patologi e popolazione è molto inferiore rispetto ai paesi ad alto reddito, la diagnosi autonoma tramite IA potrebbe in linea di principio fornire un livello di servizio diagnostico che altrimenti sarebbe irraggiungibile \cite{elendu2023}.

Tuttavia, il concetto di diagnosi autonoma tramite IA affronta sostanziali barriere normative, etiche e pratiche. I quadri normativi nella maggior parte delle giurisdizioni richiedono che i referti diagnostici siano firmati da un medico qualificato che si assume la responsabilità professionale della diagnosi. Un sistema di IA autonomo non può assumersi la responsabilità professionale in questo senso; non può essere ritenuto responsabile, non può essere disciplinato e non può essere citato in giudizio per negligenza. L'approvazione normativa dei sistemi diagnostici autonomi di IA richiederebbe quindi o una revisione fondamentale del concetto di responsabilità professionale in medicina o lo sviluppo di nuovi quadri giuridici per la responsabilità dell'IA che non esistono ancora. Inoltre, le prestazioni tecniche degli attuali sistemi di IA — sebbene impressionanti in compiti ben definiti — non sono sufficienti a giustificare la rimozione della supervisione umana nell'intera gamma di scenari diagnostici incontrati nella pratica patologica di routine \cite{zhang2023}.

\subsection{Barriere Normative ed Etiche alla Piena Automazione}

Le barriere etiche alla piena automazione in patologia sono altrettanto significative di quelle normative. La rimozione della supervisione umana dai processi diagnostici solleva profonde domande sui valori incorporati nella pratica medica. La medicina non è semplicemente un'impresa tecnica; è un'impresa morale, in cui i professionisti hanno obblighi verso i pazienti che vanno oltre l'identificazione accurata della malattia. Questi obblighi includono il dovere di comunicare le diagnosi con compassione e sensibilità, di considerare i valori e le preferenze del paziente nel contesto dell'incertezza diagnostica e di esercitare il giudizio professionale nei casi che esulano dall'ambito dei protocolli stabiliti. Queste dimensioni della pratica medica non possono essere automatizzate, e la loro importanza non dovrebbe essere diminuita dall'entusiasmo per l'efficienza tecnica \cite{hanna2024}.

Da una prospettiva pratica, l'implementazione di sistemi di IA autonomi in patologia richiederebbe un livello di affidabilità tecnica — in termini di accuratezza, robustezza allo spostamento del dominio e resistenza agli input avversariali — che i sistemi attuali non raggiungono. Richiederebbe anche meccanismi robusti per rilevare e rispondere ai guasti del sistema, per gestire i casi che esulano dalla competenza del modello e per mantenere la qualità dei dati di addestramento da cui il sistema dipende. Lo sviluppo di questi meccanismi è un'area di ricerca attiva, ma è lontano dall'essere completo. Nel frattempo, il modello appropriato per l'IA in patologia è quello della \emph{collaborazione uomo-IA}, in cui gli strumenti di IA potenziano e supportano l'expertise umana piuttosto che sostituirla, e in cui la supervisione umana viene mantenuta come salvaguardia contro gli errori dell'IA \cite{song2023}.

\subsection{Il Futuro dell'Automazione di Laboratorio}

Guardando più avanti, l'integrazione di robotica, IA e automazione di laboratorio è destinata a trasformare l'ambiente fisico del laboratorio di patologia in modi difficili da prevedere in dettaglio ma i cui contorni generali stanno diventando visibili. I laboratori del futuro potrebbero essere caratterizzati da un funzionamento continuo, 24 ore su 24, abilitato da sistemi robotici che non richiedono riposo; da un monitoraggio della qualità in tempo reale che rileva i guasti di processazione prima che influenzino la qualità diagnostica; da un'integrazione senza soluzione di continuità di dati istologici, molecolari e di imaging in piattaforme diagnostiche unificate; e da sistemi di IA in grado di triaggiare i casi per urgenza, instradare i vetrini agli specialisti appropriati e segnalare i casi per la discussione multidisciplinare. Questi sviluppi richiederanno ai tecnici di patologia e ai biologi di sviluppare nuove competenze e di assumere nuovi ruoli nella gestione e supervisione dei sistemi automatizzati \cite{elendu2023}.

% ============================================================
\section{Il Futuro Ruolo del Tecnico di Patologia}
\label{sec:ch5_technician}
% ============================================================

\subsection{L'IA come Trasformazione, non Eliminazione}

Un'ansia comune tra i tecnici di patologia e i biologi è che l'IA e l'automazione renderanno i loro ruoli obsoleti. Questa ansia è comprensibile ma, sulla base delle prove disponibili, esagerata. La storia dell'automazione di laboratorio fornisce precedenti istruttivi: l'introduzione degli analizzatori ematologici automatizzati non ha eliminato il ruolo del tecnico di ematologia ma lo ha trasformato, spostando l'attenzione dal conteggio manuale delle cellule alla gestione degli strumenti, al controllo della qualità e alla revisione dei risultati anomali segnalati. Allo stesso modo, l'introduzione dell'IA nella patologia digitale è destinata a trasformare il ruolo del tecnico piuttosto che eliminarlo, spostando l'attenzione dai compiti routinari e ripetitivi a funzioni di livello superiore che richiedono giudizio umano, comprensione contestuale e responsabilità professionale \cite{song2023}.

I compiti più probabilmente automatizzati nel breve termine sono quelli ben definiti, ripetitivi e suscettibili di formalizzazione algoritmica: screening primario dei vetrini per la presenza di tessuto, controllo della qualità della colorazione e della scansione, quantificazione dei biomarcatori standard e triage dei casi per urgenza. I compiti meno probabilmente automatizzati — e quindi più probabilmente destinati a rimanere di competenza dei professionisti umani — sono quelli che richiedono giudizio contestuale, comunicazione con i colleghi clinici, gestione di risultati inattesi e supervisione dei sistemi automatizzati. Il tecnico del futuro dovrà essere competente sia nelle competenze tradizionali della tecnica istologica che nelle nuove competenze della patologia digitale e della gestione dell'IA \cite{bankhead2022}.

\subsection{Nuove Competenze per l'Era della Patologia Digitale}

L'integrazione dell'IA nei flussi di lavoro di patologia crea una domanda di nuove competenze tecniche tra i tecnici di patologia e i biologi. L'\emph{alfabetizzazione ai dati} — la capacità di comprendere, interpretare e valutare criticamente i dati e i sistemi che li generano — è sempre più essenziale. I tecnici che lavorano con le piattaforme di patologia digitale devono comprendere i principi dell'acquisizione delle immagini e del controllo della qualità, i fattori che influenzano la qualità delle immagini e le implicazioni della qualità delle immagini per le prestazioni dell'IA. Devono essere in grado di identificare gli artefatti — pieghe del tessuto, bolle d'aria, regioni fuori fuoco, irregolarità di colorazione — che possono compromettere l'analisi dell'IA e di intraprendere azioni correttive quando tali artefatti vengono rilevati \cite{bankhead2022}.

La \emph{validazione dell'IA} è un'altra competenza critica. I tecnici responsabili del funzionamento dei sistemi diagnostici assistiti dall'IA devono comprendere i principi della validazione del modello, le metriche utilizzate per valutare le prestazioni dell'IA e le procedure per rilevare e segnalare i guasti del modello. Devono essere in grado di interpretare criticamente gli output dell'IA — di riconoscere quando una raccomandazione dell'IA è plausibile e quando dovrebbe essere messa in discussione — e di escalare le preoccupazioni ai patologi o agli amministratori di sistema quando appropriato. Le competenze di base di \emph{scripting} — la capacità di scrivere semplici script in Python o R per automatizzare le attività di gestione dei dati, generare report di controllo della qualità o eseguire analisi statistiche di base — sono sempre più apprezzate nei laboratori di patologia digitale e rappresentano un significativo differenziatore di carriera per i tecnici che le acquisiscono \cite{song2023}.

\subsection{Sviluppo Professionale Continuo e Percorsi di Carriera}

Il rapido ritmo del cambiamento nella patologia digitale e nell'IA crea sia una sfida che un'opportunità per lo sviluppo professionale continuo (CPD). La sfida è che le conoscenze e le competenze necessarie per lavorare efficacemente in un ambiente di patologia digitale si stanno evolvendo più velocemente di quanto i tradizionali quadri CPD possano accogliere; un tecnico che si è qualificato cinque anni fa potrebbe scoprire che aspetti significativi della sua formazione sono già obsoleti. L'opportunità è che lo sviluppo di nuove competenze nella patologia digitale e nell'IA apre nuovi percorsi di carriera che non erano disponibili per le generazioni precedenti di tecnici di patologia \cite{hanna2024}.

I percorsi di carriera nella patologia digitale includono ruoli nell'implementazione e gestione della patologia digitale, nella validazione dell'IA e nell'assicurazione della qualità, nell'informatica di laboratorio e nel supporto alla ricerca. Alcuni tecnici possono progredire verso ruoli di specialisti di patologia digitale o coordinatori dell'IA, assumendo la responsabilità della gestione delle piattaforme di patologia digitale e dei sistemi di IA in un laboratorio o in una rete ospedaliera. Altri possono passare a ruoli industriali con aziende che sviluppano hardware, software o strumenti di IA per la patologia digitale. Gli organismi professionali, incluso l'Institute of Biomedical Science (IBMS) nel Regno Unito, hanno iniziato a sviluppare quadri CPD e qualifiche specialistiche nella patologia digitale che forniscono un riconoscimento formale di queste nuove competenze. Il coinvolgimento con questi quadri è fortemente incoraggiato per i tecnici che desiderano sviluppare la propria carriera in questo settore.

\subsection{L'Importanza della Supervisione Umana}

Qualunque sia la futura traiettoria dell'IA e dell'automazione in patologia, l'importanza della supervisione umana non può essere sopravvalutata. I sistemi di IA possono fallire in modi inaspettati, sottili e difficili da rilevare senza un monitoraggio attivo. Le conseguenze dei guasti dell'IA non rilevati in un contesto diagnostico possono essere gravi — diagnosi mancate, decisioni di trattamento errate, danni ai pazienti. La supervisione umana — il coinvolgimento attivo e informato dei professionisti qualificati con gli output dell'IA — è la principale salvaguardia contro queste conseguenze. Questa supervisione non è passiva; richiede professionisti che comprendano le capacità e le limitazioni dei sistemi di IA con cui lavorano, che mantengano le competenze necessarie per rilevare gli errori dell'IA e che siano autorizzati a mettere in discussione gli output dell'IA quando il loro giudizio professionale indica che qualcosa non va \cite{zhang2023}.

Il concetto di controllo umano significativo, introdotto nella Sezione~\ref{sec:ch5_philosophy}, è direttamente rilevante qui. Affinché la supervisione umana sia significativa, deve essere più di un esercizio di timbro di gomma in cui un professionista approva formalmente gli output dell'IA senza confrontarsi genuinamente con essi. Richiede che i professionisti abbiano tempo, informazioni e competenze sufficienti per valutare criticamente gli output dell'IA, e che le culture istituzionali e i flussi di lavoro supportino piuttosto che minino questo impegno critico. Creare le condizioni per una supervisione umana significativa è una responsabilità condivisa degli sviluppatori di IA, delle istituzioni che li implementano, degli organismi professionali e dei regolatori — ed è una responsabilità che il tecnico di patologia del futuro sarà chiamato a esercitare ogni giorno \cite{elendu2023}.

% ============================================================
%  AI Ethics Framework Figure
% ============================================================

\begin{figure}[htbp]
    \centering
    \usetikzlibrary{arrows.meta, positioning, calc}

\begin{tikzpicture}[
    every node/.style={font=\small\sffamily},
    main/.style={
        rectangle, rounded corners=5pt,
        fill={rgb,255:red,46;green,134;blue,171},
        text=white,
        minimum width=3.2cm, minimum height=0.85cm,
        align=center, font=\small\sffamily\bfseries
    },
    ann/.style={
        rectangle, rounded corners=3pt,
        fill={rgb,255:red,255;green,183;blue,77},
        text={rgb,255:red,70;green,45;blue,0},
        minimum width=2.9cm, minimum height=0.62cm,
        align=center, font=\footnotesize\sffamily,
        draw={rgb,255:red,220;green,145;blue,25},
        line width=0.5pt
    },
    arr/.style={
        -{Stealth[length=5pt,width=4pt]},
        line width=0.9pt,
        color={rgb,255:red,46;green,134;blue,171}
    },
    dann/.style={
        dashed,
        -{Stealth[length=4pt,width=3pt]},
        line width=0.5pt,
        color={rgb,255:red,180;green,120;blue,30}
    },
    stg/.style={
        font=\footnotesize\sffamily\itshape,
        text={rgb,255:red,46;green,134;blue,171}
    }
]

%% ── Stage labels ────────────────────────────────────────────────────
\node[stg] at (0,   0.72) {Stage 1};
\node[stg] at (4.2, 0.72) {Stage 2};
\node[stg] at (8.4, 0.72) {Stage 3};

%% ── Main pipeline ───────────────────────────────────────────────────
\node[main] (dc) at (0,   0) {Data\\Collection};
\node[main] (md) at (4.2, 0) {Model\\Development};
\node[main] (cd) at (8.4, 0) {Clinical\\Deployment};

\draw[arr] (dc) -- (md);
\draw[arr] (md) -- (cd);

%% ── Annotation boxes: Data Collection ──────────────────────────────
\node[ann] (dc1) at (0.0, -1.55) {Bias in training data};
\node[ann] (dc2) at (0.0, -2.28) {Privacy \& consent};
\node[ann] (dc3) at (0.0, -3.01) {Representativeness};

\draw[dann] (dc.south)  -- (dc1.north);
\draw[dann] (dc1.south) -- (dc2.north);
\draw[dann] (dc2.south) -- (dc3.north);

%% ── Annotation boxes: Model Development ────────────────────────────
\node[ann] (md1) at (4.2, -1.55) {Algorithmic fairness};
\node[ann] (md2) at (4.2, -2.28) {Interpretability};
\node[ann] (md3) at (4.2, -3.01) {Validation};

\draw[dann] (md.south)  -- (md1.north);
\draw[dann] (md1.south) -- (md2.north);
\draw[dann] (md2.south) -- (md3.north);

%% ── Annotation boxes: Clinical Deployment ───────────────────────────
\node[ann] (cd1) at (8.4, -1.55) {Accountability};
\node[ann] (cd2) at (8.4, -2.28) {Human oversight};
\node[ann] (cd3) at (8.4, -3.01) {Regulatory compliance};

\draw[dann] (cd.south)  -- (cd1.north);
\draw[dann] (cd1.south) -- (cd2.north);
\draw[dann] (cd2.south) -- (cd3.north);

%% ── Bottom rule ─────────────────────────────────────────────────────
\draw[gray!30, line width=0.4pt]
    (-1.8, -3.42) -- (10.2, -3.42);

\end{tikzpicture}

    \caption{Il quadro etico dell'IA in patologia, che illustra la pipeline dalla raccolta dei dati attraverso lo sviluppo del modello fino all'implementazione clinica. In ogni fase emergono considerazioni etiche fondamentali: la raccolta dei dati è soggetta al \emph{bias dei dati} (sottorappresentazione dei gruppi demografici, variazione degli scanner ed effetti batch); lo sviluppo del modello deve affrontare l'\emph{equità} (prestazioni eque tra i sottogruppi) e la \emph{spiegabilità} (output interpretabili tramite Grad-CAM, mappe di attenzione e valori SHAP); e l'implementazione clinica richiede la \emph{responsabilità} (linee chiare di responsabilità per gli errori dell'IA), la \emph{conformità normativa} (Regolamento UE sull'IA, quadro FDA SaMD) e la continua \emph{sorveglianza post-commercializzazione}. La supervisione umana è rappresentata come un requisito trasversale che abbraccia tutte le fasi della pipeline.}
    \label{fig:ethics_framework}
\end{figure}

% ============================================================
\section{Attività Pratiche}
\label{sec:ch5_activities}
% ============================================================

\subsection{Attività 1: Dibattito di Gruppo --- Fiducia e Responsabilità nella Patologia Assistita dall'IA}

\textbf{Obiettivo:} Sviluppare la capacità degli studenti di ragionamento etico critico sull'IA in patologia ed esplorare le tensioni tra i potenziali benefici dell'IA e i rischi di bias, errori e lacune di responsabilità.

\textbf{Formato:} Dibattito strutturato in gruppi da sei a otto studenti, con due squadre che sostengono posizioni opposte su una mozione. Le mozioni suggerite includono: \emph{``Questa assemblea ritiene che i sistemi di IA debbano essere autorizzati a emettere referti diagnostici senza revisione umana obbligatoria in contesti con risorse limitate''} o \emph{``Questa assemblea ritiene che i rischi del bias algoritmico nella patologia assistita dall'IA superino i potenziali benefici per le popolazioni svantaggiate.''}

\textbf{Preparazione:} Ogni squadra dovrebbe prepararsi per 30 minuti, attingendo al materiale di questo capitolo e ad almeno due fonti di letteratura primaria. Le squadre dovrebbero considerare le prospettive di pazienti, patologi, tecnici, regolatori e sviluppatori di IA. Ogni squadra dovrebbe designare un oratore principale, un oratore di replica e un verbalizzatore.

\textbf{Procedura:} Il dibattito dovrebbe seguire un formato strutturato: dichiarazioni di apertura (5 minuti per squadra), replica (3 minuti per squadra), domande aperte (10 minuti) e dichiarazioni di chiusura (2 minuti per squadra). Un facilitatore dovrebbe garantire che tutti i partecipanti abbiano l'opportunità di contribuire e dovrebbe sollecitare le squadre a considerare prospettive che potrebbero aver trascurato.

\textbf{Debriefing:} Dopo il dibattito, il facilitatore dovrebbe condurre una discussione di debriefing di 15 minuti esplorando le seguenti domande: (1) Quali argomenti sono stati più persuasivi e perché? (2) C'erano prospettive che nessuna delle due squadre ha adeguatamente rappresentato? (3) Come ha cambiato il dibattito il vostro pensiero sulle dimensioni etiche dell'IA in patologia? (4) Quali meccanismi istituzionali o normativi sarebbero necessari per affrontare le preoccupazioni sollevate? Gli studenti dovrebbero presentare una riflessione individuale di 500 parole sul dibattito entro una settimana.

\textbf{Risultati di apprendimento:} Al completamento di questa attività, gli studenti saranno in grado di articolare i principali argomenti etici a favore e contro la diagnosi autonoma tramite IA; identificare le prospettive degli stakeholder rilevanti per le decisioni di implementazione dell'IA; e applicare quadri etici a scenari nuovi nella patologia assistita dall'IA.

\subsection{Attività 2: Workshop --- Carriere e Tendenze Future nella Patologia Digitale}

\textbf{Obiettivo:} Aiutare gli studenti a comprendere il panorama in evoluzione delle opportunità di carriera nella patologia digitale e nell'IA, e a sviluppare un piano personale di sviluppo professionale che incorpori competenze digitali e di IA.

\textbf{Formato:} Workshop di mezza giornata che combina una discussione a panel con relatori invitati, esercizi in piccoli gruppi e riflessione individuale.

\textbf{Discussione a panel (60 minuti):} Invitare da tre a quattro relatori che rappresentano diversi percorsi di carriera nella patologia digitale: un patologo consulente con un interesse subspecialistico nella patologia digitale; un biologo o tecnico di patologia che lavora in un ruolo di implementazione della patologia digitale; un data scientist o ingegnere dell'IA che lavora in un'azienda di IA per la patologia; e, se possibile, un difensore dei pazienti o un rappresentante di un'organizzazione di pazienti con interesse nell'IA in sanità. Ogni relatore dovrebbe tenere una presentazione di 10 minuti sul proprio percorso di carriera e sulla propria prospettiva sul futuro della patologia digitale, seguita da 20 minuti di domande degli studenti.

\textbf{Esercizio in piccoli gruppi (45 minuti):} Gli studenti lavorano in gruppi da quattro a cinque per mappare le competenze richieste per tre ruoli di carriera nella patologia digitale di loro scelta, utilizzando un modello strutturato che copre: (a) competenze tecniche (analisi delle immagini, scripting, validazione dell'IA); (b) conoscenze scientifiche (istologia, patologia molecolare, statistica); (c) competenze professionali (comunicazione, gestione della qualità, consapevolezza normativa); e (d) percorsi di sviluppo professionale continuo (qualifiche IBMS, corsi online, iscrizione a società professionali). I gruppi presentano i loro risultati all'intero gruppo in un riassunto di cinque minuti.

\textbf{Riflessione individuale (30 minuti):} Ogni studente completa un piano personale di sviluppo professionale identificando: (a) due competenze di patologia digitale o IA che desidera sviluppare nei prossimi 12 mesi; (b) le azioni specifiche che intraprenderà per sviluppare queste competenze (corsi, letture, esperienza pratica); e (c) come dimostrerà queste competenze a un futuro datore di lavoro o organismo professionale. Gli studenti sono incoraggiati a riesaminare e aggiornare questo piano a intervalli di sei mesi durante la loro formazione post-laurea.

\textbf{Risultati di apprendimento:} Al completamento di questa attività, gli studenti saranno in grado di descrivere almeno tre percorsi di carriera nella patologia digitale e nell'IA; identificare le competenze tecniche, scientifiche e professionali richieste per i ruoli nella patologia digitale; e costruire un piano personale di sviluppo professionale che incorpori competenze digitali e di IA.

% ============================================================
\section{Sintesi}
\label{sec:ch5_summary}
% ============================================================

Questo capitolo ha esaminato le sfide, le dimensioni etiche e le prospettive future della patologia assistita dall'IA da molteplici angolazioni. Le limitazioni tecniche degli attuali sistemi di IA — inclusi il fallimento della generalizzazione, il bias del dataset, l'illusione dell'oggettività, l'obsolescenza del modello, il divario tra prova di concetto e implementazione clinica e i costi dell'infrastruttura computazionale — ci ricordano che gli strumenti di IA sono potenti ma imperfetti, e che la loro implementazione richiede una validazione rigorosa, un monitoraggio continuo e un impegno critico da parte dei professionisti \cite{bankhead2022,howard2021}.

L'analisi del bias nei sistemi di IA — che comprende il bias dei dati, il bias delle etichette, il bias di misurazione e il bias di implementazione — evidenzia il rischio che gli strumenti di IA esacerbino le disuguaglianze sanitarie esistenti se vengono sviluppati e implementati senza un'attenzione esplicita all'equità \cite{chen2023,seyyedkalantari2021}. Il campo dell'IA spiegabile offre soluzioni parziali ma imperfette al problema della scatola nera, fornendo strumenti come Grad-CAM, mappe di attenzione e valori SHAP che possono rendere gli output dell'IA più interpretabili, mentre le linee guida SPIRIT-AI e CONSORT-AI forniscono quadri per la segnalazione trasparente della ricerca sull'IA \cite{amann2020,cruzirivera2020}.

La riflessione filosofica sulla natura della diagnosi automatica, il ruolo della conoscenza tacita in patologia e lo status epistemologico della diagnosi assistita dall'IA rivela che l'integrazione dell'IA in patologia non è semplicemente un evento tecnico ma una trasformazione del panorama epistemico e morale della medicina diagnostica \cite{hanna2024,song2023}. Le dimensioni etiche della diagnosi assistita dall'IA — che comprendono il consenso informato, la privacy dei dati, la proprietà intellettuale, la responsabilità e la conformità normativa — richiedono lo sviluppo di nuovi quadri di governance e nuove competenze professionali \cite{zhang2023,elendu2023}.

La discussione sulla robotica, l'automazione e i sistemi autonomi illustra la traiettoria verso flussi di lavoro di patologia sempre più automatizzati, evidenziando al contempo le barriere normative ed etiche alla piena automazione e l'importanza duratura della supervisione umana. Infine, l'analisi del futuro ruolo del tecnico di patologia sottolinea che l'IA e l'automazione trasformeranno piuttosto che elimineranno il ruolo del tecnico, creando nuove esigenze di alfabetizzazione ai dati, competenze di validazione dell'IA e competenze di patologia digitale, e aprendo nuovi percorsi di carriera per coloro che abbracciano questi cambiamenti \cite{bankhead2022,song2023}.

Il messaggio generale di questo capitolo è che l'integrazione responsabile dell'IA in patologia richiede non solo eccellenza tecnica ma anche consapevolezza etica, pensiero critico e un impegno per i valori di equità, trasparenza e responsabilità. Questi non sono optional; sono competenze fondamentali per il professionista della patologia del ventunesimo secolo.

% ============================================================
\section{Domande di Revisione}
\label{sec:ch5_questions}
% ============================================================

\begin{enumerate}

    \item \label{ch5:q1} Un modello di apprendimento profondo addestrato su immagini di vetrino intero di una singola istituzione raggiunge il 95\% di accuratezza su un set di test trattenuto della stessa istituzione, ma solo il 72\% di accuratezza quando applicato a vetrini di un'istituzione diversa che utilizza uno scanner diverso. Questo fenomeno è meglio descritto come:
    \begin{enumerate}[label=\Alph*.]
        \item Bias delle etichette derivante dalla variabilità inter-osservatore nelle annotazioni
        \item Spostamento del dominio dovuto a differenze nelle condizioni di acquisizione delle immagini
        \item Deriva del concetto causata da cambiamenti nei criteri diagnostici nel tempo
        \item Bias di implementazione risultante da differenze nei dati demografici dei pazienti
    \end{enumerate}

    \item \label{ch5:q2} Howard e colleghi hanno dimostrato che i modelli addestrati sul dataset The Cancer Genome Atlas (TCGA) possono inavvertitamente imparare a classificare i vetrini per istituzione contribuente piuttosto che per biologia tumorale. La causa principale di questo fenomeno è:
    \begin{enumerate}[label=\Alph*.]
        \item Capacità insufficiente del modello per apprendere caratteristiche specifiche del tumore
        \item Firme specifiche del sito derivanti da differenze nella processazione del tessuto e nella scansione tra le istituzioni contribuenti
        \item Etichettatura deliberatamente errata dei vetrini da parte delle istituzioni contribuenti
        \item L'uso del transfer learning da dataset di immagini non patologiche
    \end{enumerate}

    \item \label{ch5:q3} Quale delle seguenti descrive meglio il concetto di ``bias da automazione'' nel contesto della patologia assistita dall'IA?
    \begin{enumerate}[label=\Alph*.]
        \item La tendenza dei modelli di IA a ottenere prestazioni migliori nei compiti automatizzati rispetto ai compiti che richiedono il giudizio umano
        \item Le prestazioni sistematicamente inferiori dei modelli di IA per i pazienti di gruppi demografici sottorappresentati
        \item La tendenza dei professionisti umani ad accettare acriticamente le raccomandazioni dell'IA senza una valutazione indipendente
        \item Il bias introdotto nei modelli di IA dall'uso di strumenti di annotazione automatizzati
    \end{enumerate}

    \item \label{ch5:q4} La Mappatura dell'Attivazione di Classe Ponderata dal Gradiente (Grad-CAM) è un metodo utilizzato nell'IA spiegabile. Quale delle seguenti descrive più accuratamente come Grad-CAM genera il suo output?
    \begin{enumerate}[label=\Alph*.]
        \item Riaddestra il modello su un sottoinsieme dei dati di addestramento per identificare gli esempi di addestramento più influenti
        \item Calcola il gradiente dell'output del modello rispetto alle attivazioni di uno strato convoluzionale per produrre una mappa di calore spaziale
        \item Utilizza i valori di Shapley game-teorici per attribuire l'output del modello alle singole caratteristiche di input
        \item Genera immagini controfattuali perturbando l'input e osservando i cambiamenti nell'output del modello
    \end{enumerate}

    \item \label{ch5:q5} Seyyed-Kalantari e colleghi hanno dimostrato che i modelli di IA addestrati su grandi dataset pubblici hanno sistematicamente sottodiagnosticato le condizioni nei pazienti di gruppi emarginati. Questa scoperta è più direttamente rilevante per quale tipo di bias dell'IA?
    \begin{enumerate}[label=\Alph*.]
        \item Bias di misurazione derivante dalla variazione degli scanner
        \item Bias delle etichette derivante dalla variabilità inter-osservatore
        \item Bias dei dati derivante dalla sottorappresentazione dei gruppi demografici nei dataset di addestramento
        \item Bias di implementazione derivante da differenze nel flusso di lavoro clinico
    \end{enumerate}

    \item \label{ch5:q6} Le linee guida di reporting CONSORT-AI sono state sviluppate per affrontare quale delle seguenti lacune nella segnalazione della ricerca sull'IA?
    \begin{enumerate}[label=\Alph*.]
        \item L'assenza di metodi standardizzati per l'addestramento di modelli di apprendimento profondo su dataset di patologia
        \item La mancanza di requisiti per la divulgazione di informazioni specifiche dell'IA nei report degli studi clinici che coinvolgono interventi di IA
        \item Il fallimento delle linee guida esistenti nell'affrontare i requisiti computazionali dell'addestramento dei modelli di IA
        \item L'assenza di quadri normativi per l'approvazione dei dispositivi medici basati sull'IA
    \end{enumerate}

    \item \label{ch5:q7} Nel contesto della patologia assistita dall'IA, il concetto di ``controllo umano significativo'' richiede quale delle seguenti?
    \begin{enumerate}[label=\Alph*.]
        \item Che un professionista umano sia fisicamente presente in laboratorio quando i sistemi di IA sono in funzione
        \item Che i sistemi di IA siano progettati per replicare i processi decisionali umani nel modo più fedele possibile
        \item Che i professionisti umani abbiano una comprensione sufficiente degli output dell'IA per esercitare una supervisione genuina e informata delle decisioni diagnostiche finali
        \item Che i sistemi di IA vengano riaddestrati da esperti umani a intervalli regolari per prevenire la deriva del concetto
    \end{enumerate}

    \item \label{ch5:q8} Il Regolamento UE sull'IA classifica i sistemi di IA utilizzati a fini diagnostici in sanità in quale categoria di rischio, e qual è la principale conseguenza di questa classificazione?
    \begin{enumerate}[label=\Alph*.]
        \item Rischio basso; i sistemi diagnostici di IA sono esenti dai requisiti normativi ma devono mostrare un avviso di trasparenza
        \item Rischio moderato; i sistemi diagnostici di IA devono essere registrati presso le autorità nazionali ma non richiedono una valutazione di conformità
        \item Rischio alto; i sistemi diagnostici di IA sono soggetti a requisiti di trasparenza, supervisione umana, robustezza e valutazione di conformità prima dell'implementazione
        \item Rischio inaccettabile; i sistemi diagnostici di IA sono vietati dall'uso in contesti clinici senza esenzione speciale
    \end{enumerate}

    \item \label{ch5:q9} Quale delle seguenti descrive meglio il concetto di ``deriva del concetto'' nel contesto dei modelli di IA implementati in patologia?
    \begin{enumerate}[label=\Alph*.]
        \item Il graduale miglioramento delle prestazioni del modello che si verifica quando il modello viene esposto a più dati dopo l'implementazione
        \item Il degrado delle prestazioni del modello che si verifica quando le relazioni statistiche nell'ambiente di implementazione cambiano rispetto a quelle nell'ambiente di addestramento
        \item La tendenza dei modelli di IA a generare previsioni sempre più sicure nel tempo, indipendentemente dalla qualità dell'input
        \item Il cambiamento nell'architettura del modello che si verifica quando un modello viene aggiornato per incorporare nuovi dati di addestramento
    \end{enumerate}

    \item \label{ch5:q10} Un laboratorio di patologia sta valutando l'implementazione di un sistema di IA per lo screening primario dei vetrini di citologia cervicale. Quale delle seguenti combinazioni di azioni rifletterebbe meglio i principi di implementazione responsabile dell'IA discussi in questo capitolo?
    \begin{enumerate}[label=\Alph*.]
        \item Implementare il sistema immediatamente sulla base delle prestazioni benchmark pubblicate e monitorare i reclami del personale clinico
        \item Validare il sistema su un campione rappresentativo di vetrini della popolazione di pazienti locale, stabilire un protocollo di revisione umana per tutti gli output dell'IA, implementare la sorveglianza post-commercializzazione e ottenere l'approvazione normativa appropriata
        \item Fare affidamento sui dati di validazione del produttore e sull'approvazione normativa, senza condurre una validazione locale, poiché ciò duplicherebbe le prove esistenti
        \item Implementare il sistema in modalità completamente autonoma per massimizzare il throughput, con revisione umana riservata ai casi segnalati come incerti dall'IA
    \end{enumerate}

\end{enumerate}
